\documentclass{article}
\usepackage[margin=1in]{geometry}
\usepackage{amsmath,amssymb}
\usepackage[most]{tcolorbox}
\usepackage[hidelinks]{hyperref}

\newcommand{\BookName}{Pattern Recognition and Machine Learning}
\newcommand{\BookAuthor}{Christopher M. Bishop}
\newcommand{\ChapterName}{Introduction}
\newcommand{\ExerciseID}{1.11}

\title{Chapter: \ChapterName}
\author{Ahonakpon Guy Gbaguidi}
\date{}

\begin{document}

\maketitle

\begin{center}
  \large\textbf{Exercise \ExerciseID}

  \vspace{0.5em}

  \normalsize\BookName\ - \BookAuthor
\end{center}

\begin{tcolorbox}[breakable,colback=gray!5,colframe=gray!60,title=Solution]
By setting the derivative of the log likelihood function (1.54) with respect to 
\(\mu\) and \(\sigma^2\) equal to zero.

Let restate the equation (1.54) for clarity:

\begin{equation}
  \ln p(\mathbf{x}|\mu,\sigma^2) = - \frac{1}{2\sigma^2}\sum_{n=1}^N (x_n - \mu)^2 -\frac{N}{2}\ln(\sigma^2) - \frac{N}{2}\ln(2\pi) \tag{1.54}
\end{equation}

In this exercise, we want to verify the results (1.55) and (1.56).

\paragraph{Verification of (1.55).}

\begin{equation}
  \mu_{\text{ML}} = \frac{1}{N} \sum_{n=1}^N x_n \tag{1.55}
\end{equation}

Let start by computing the derivative of the log likelihood function with respect to \(\mu\):

\begin{equation}
  \frac{\partial}{\partial \mu} \ln p(\mathbf{x}|\mu,\sigma^2) = -\frac{1}{2\sigma^2} \sum_{n=1}^N \frac{\partial}{\partial \mu} (x_n - \mu)^2
\end{equation}

because the other terms do not depend on \(\mu\). Now, we can compute the derivative inside the sum:

\begin{equation}
  \frac{\partial}{\partial \mu} (x_n - \mu)^2 = -2(x_n - \mu)
\end{equation}

Substituting this back into the expression for the derivative of the log likelihood, we get:

\begin{equation}
  \frac{\partial}{\partial \mu} \ln p(\mathbf{x}|\mu,\sigma^2) = -\frac{1}{2\sigma^2} \sum_{n=1}^N (-2)(x_n - \mu) = \frac{1}{\sigma^2} \sum_{n=1}^N (x_n - \mu)
\end{equation}

We can now set this derivative equal to zero to find the maximum likelihood estimate for \(\mu\):

\begin{align*}
  \frac{\partial}{\partial \mu} \ln p(\mathbf{x}|\mu,\sigma^2) &= 0 \\
  \frac{1}{\sigma^2} \sum_{n=1}^N (x_n - \mu) &= 0 \\
  \sum_{n=1}^N (x_n - \mu) &= 0 \quad \text{(since } \sigma^2 > 0) \\
  \sum_{n=1}^N x_n - N\mu &= 0 \\
  N\mu &= \sum_{n=1}^N x_n \\
  \mu &= \frac{1}{N} \sum_{n=1}^N x_n
\end{align*}

This matches the result given in (1.55).

\paragraph{Verification of (1.56).}
\begin{equation}
  \sigma^2_{\text{ML}} = \frac{1}{N} \sum_{n=1}^N (x_n - \mu_{\text{ML}})^2 \tag{1.56}
\end{equation}

Now, let's compute the derivative of the log likelihood function with respect to \(\sigma^2\):

\begin{equation}
  \frac{\partial}{\partial \sigma^2} \ln p(\mathbf{x}|\mu,\sigma^2) = -\frac{1}{2} \sum_{n=1}^N (x_n - \mu)^2 \cdot \frac{\partial}{\partial \sigma^2} \left( \frac{1}{\sigma^2} \right) - \frac{N}{2} \cdot \frac{\partial}{\partial \sigma^2} \ln(\sigma^2)
\end{equation}

We can compute the derivatives inside the expression:

\begin{align*}
  \frac{\partial}{\partial \sigma^2} \left( \frac{1}{\sigma^2} \right) &= -\frac{1}{\sigma^4} \\
  \frac{\partial}{\partial \sigma^2} \ln(\sigma^2) &= \frac{1}{\sigma^2}
\end{align*}

Substituting these back into the expression for the derivative of the log likelihood, we get:

\begin{equation}
  \frac{\partial}{\partial \sigma^2} \ln p(\mathbf{x}|\mu,\sigma^2) = -\frac{1}{2} \sum_{n=1}^N (x_n - \mu)^2 \cdot \left( -\frac{1}{\sigma^4} \right) - \frac{N}{2} \cdot \frac{1}{\sigma^2}
\end{equation}

Simplifying this expression:

\begin{equation}
  \frac{\partial}{\partial \sigma^2} \ln p(\mathbf{x}|\mu,\sigma^2) = \frac{1}{2\sigma^4} \sum_{n=1}^N (x_n - \mu)^2 - \frac{N}{2\sigma^2}
\end{equation}

We can now set this derivative equal to zero to find the maximum likelihood estimate for \(\sigma^2\):

\begin{align*}
  \frac{\partial}{\partial \sigma^2} \ln p(\mathbf{x}|\mu,\sigma^2) &= 0 \\
  \frac{1}{2\sigma^4} \sum_{n=1}^N (x_n - \mu)^2 - \frac{N}{2\sigma^2} &= 0 \\
  \sum_{n=1}^N (x_n - \mu)^2 &= N\sigma^2 \\
  \sigma^2 &= \frac{1}{N} \sum_{n=1}^N (x_n - \mu)^2
\end{align*}

This matches the result given in (1.56).

\vspace{1em}

\textbf{ We can now conclude that the maximum likelihood estimates for \(\mu\) and \(\sigma^2\) are 
given by (1.55) and (1.56) respectively, and both estimates are unique because the log likelihood function is concave with respect to \(\mu\) and \(\sigma^2\).}
\end{tcolorbox}

\end{document}